% TYK §2.1.9 + Ek-8: Simgeler alfabetik sıralanır ve **birimleri belirtilir**.
% Format: "simge : Açıklama, (birim)". Birimsiz (boyutsuz) niceliklerde birim
% yerine "(boyutsuz)" yazılır.
%
% Sıralama: Latin harfler (alfabetik), ardından Yunan harfleri.
\begin{simge}
	\item [$F_1$]          F1 skoru, (boyutsuz)
	\item [$h$]            Tahmin ufku, (saat)
	\item [$L$]            Kayıp fonksiyonu değeri, (boyutsuz)
	\item [$n$]            Gözlem sayısı, (adet)
	\item [$P$]            Kesinlik (precision), (boyutsuz)
	\item [$R$]            Duyarlılık (recall), (boyutsuz)
	\item [$R^2$]          Belirleme katsayısı, (boyutsuz)
	\item [$s$]            Log anomali skoru, (boyutsuz)
	\item [$T$]            Zaman penceresi uzunluğu, (saat)
	\item [$y$]            Gerçek kaynak kullanım değeri, (\%)
	\item [$\hat{y}$]      Tahmin edilen kaynak kullanım değeri, (\%)
	\item [$\theta$]       Kapasite alarm eşiği, (\%)
	\item [$\mu$]          Kayan ortalama, (boyutsuz)
	\item [$\phi$]         SHAP açıklama değeri, (boyutsuz)
	\item [$\sigma$]       Kayan standart sapma, (boyutsuz)
\end{simge}
